\section*{Hoofdstuk \ref{ch:g5:weather-and-climate} --- Weer en klimaat}

\begin{solution}{exo:g5:weather-and-climate:1}
(a) \omterm{def:g1:watching-the-weather:weather}{Weer} (alleen vanochtend). (b) \omterm{def:g5:weather-and-climate:climate}{Klimaat} (elke zomer). (c) \omterm{def:g1:watching-the-weather:weather}{Weer}
(gistermiddag). (d) \omterm{def:g5:weather-and-climate:climate}{Klimaat} (elke maand, jaar na jaar).
\end{solution}

\begin{solution}{exo:g5:weather-and-climate:2}
Het \omterm{def:g5:weather-and-climate:climate}{klimaat} van een plaats is het \omterm{def:g1:watching-the-weather:weather}{weer} dat die gewoonlijk heeft, maand na
maand, over veel jaren. Weerkundigen nemen het gemiddelde van de metingen
van 30 jaar.
\end{solution}

\begin{solution}{exo:g5:weather-and-climate:3}
De koudste maand is december (ongeveer 6~$^\circ$C; januari is bijna even
koud). De natste maand is november (ongeveer 160~mm). Nee: de lijn komt
nooit onder 0~$^\circ$C.
\end{solution}

\begin{solution}{exo:g5:weather-and-climate:4}
Juli is de warmste maand, met ongeveer 5~$^\circ$C. Geen enkele maand is
warmer dan 10~$^\circ$C.
\end{solution}

\begin{solution}{exo:g5:weather-and-climate:5}
De \omterm{def:g3:continents-and-oceans:equator}{evenaar}: de warme zone. De \omterm{def:g3:continents-and-oceans:poles}{Noordpool}: de koude zone (de noordelijke).
Halverwege tussen de \omterm{def:g3:continents-and-oceans:equator}{evenaar} en de \omterm{def:g3:continents-and-oceans:poles}{Zuidpool}: de zuidelijke gematigde zone
(die ligt tussen de \omterm{def:g5:weather-and-climate:tropics}{Steenbokskeerkring} en de \omterm{def:g5:weather-and-climate:tropics}{zuidpoolcirkel}).
\end{solution}

\begin{solution}{exo:g5:weather-and-climate:6}
Hilo: de warme zone. Phoenix en Seattle: de noordelijke gematigde zone.
Utqiaġvik: de noordelijke koude zone, voorbij de \omterm{def:g5:weather-and-climate:tropics}{noordpoolcirkel}.
\end{solution}

\begin{solution}{exo:g5:weather-and-climate:7}
Het hele jaar: $147 + 96 + 106 + 81 + 48 + 37 + 15 + 25 + 41 + 99 + 160 + 145
= 1\,000$ millimeter. December, januari en februari:
$145 + 147 + 96 = 388$ millimeter, meer dan een derde van de \omterm{def:g4:clouds-rain-snow:precipitation}{regen} van het
jaar in drie maanden.
\end{solution}

\begin{solution}{exo:g5:weather-and-climate:8}
Phoenix: $35 - 14 = 21$ graden. Hilo: $25 - 22 = 3$ graden. De seizoenen
veranderen de temperatuur veel meer in Phoenix; in Hilo is ze het hele jaar
bijna gelijk.
\end{solution}

\begin{solution}{exo:g5:weather-and-climate:9}
Een regenwoudklimaat: het is het hele jaar heet (27 of 28~$^\circ$C, zonder
winter), en er valt elke maand veel \omterm{def:g4:clouds-rain-snow:precipitation}{regen} (minstens 160~mm, en meer dan
2\,000~mm per jaar). De plaats ligt waarschijnlijk in de warme zone, dicht
bij de \omterm{def:g3:continents-and-oceans:equator}{evenaar}.
\end{solution}

\begin{solution}{exo:g5:weather-and-climate:10}
Eerste verklaring: het strand ligt in de warme zone, waar het het hele jaar
warm is, en de besneeuwde tuin in een gematigde of koude zone, waar januari
in de winter valt. Tweede verklaring: het strand ligt op het \omterm{def:g3:continents-and-oceans:equator}{zuidelijk
halfrond}, waar januari in de zomer valt, en de besneeuwde tuin op het
\omterm{def:g3:continents-and-oceans:equator}{noordelijk halfrond}, waar januari in de winter valt: de seizoenen zijn
omgekeerd tussen de twee helften van de \omterm{def:g3:continents-and-oceans:globe}{globe}.
\end{solution}

\begin{solution}{exo:g5:weather-and-climate:11}
Temperaturen: Phoenix is veel heter, ongeveer 35~$^\circ$C in juli tegen
ongeveer 20~$^\circ$C in augustus in Seattle, en januari in Phoenix
(ongeveer 14~$^\circ$C) is al even warm als mei in Seattle. \omterm{def:g4:clouds-rain-snow:precipitation}{Regen}: minder
dan 200~mm per jaar in Phoenix, een \omterm{def:g5:weather-and-climate:desert}{woestijn}, tegen ongeveer 1\,000~mm in
Seattle. Ze hebben dus niet hetzelfde \omterm{def:g5:weather-and-climate:climate}{klimaat}. De zones vertellen ons
ruwweg hoeveel de Zon een plaats verwarmt, van de \omterm{def:g3:continents-and-oceans:equator}{evenaar} tot de \omterm{def:g3:continents-and-oceans:poles}{polen}; ze
vertellen ons niet hoeveel \omterm{def:g4:clouds-rain-snow:precipitation}{regen} er valt, en ook niet wat het effect is van
de \omterm{def:g2:water-on-earth:ocean}{zee}, de \omterm{def:g1:landscapes-around-us:landform}{bergen} of de wind.
\end{solution}
